% Generated draft: requires mathematical and source review.
\documentclass[11pt]{article}
\usepackage[margin=25mm]{geometry}
\usepackage{fontspec}
\setmainfont{Latin Modern Roman}
\usepackage{amsmath,amssymb,mathtools,unicode-math}
\setmathfont{Latin Modern Math}
\usepackage{enumitem,needspace,xurl,hyperref}
\setlength{\parindent}{0pt}
\setlength{\parskip}{5pt}
\setlength{\emergencystretch}{3em}
\newcommand{\N}{\mathbb{N}}
\newcommand{\Z}{\mathbb{Z}}
\newcommand{\Q}{\mathbb{Q}}
\newcommand{\R}{\mathbb{R}}
\newcommand{\C}{\mathbb{C}}
\newcommand{\F}{\mathbb{F}}
\newcommand{\A}{\mathbb{A}}
\newcommand{\PP}{\mathbb{P}}
\newcommand{\E}{\mathbb{E}}
\newcommand{\eps}{\varepsilon}
\newcommand{\dd}{\,\mathrm{d}}
\begin{document}
\section*{UnsolvedMath ID: 10900049}
\textbf{Source problem code:} catalogue entry
\subsection{Definitions and mathematical statement}
Let $f\in\Z[x]$ be irreducible over $\Q$ with all roots real, let $f(\alpha)=0$, and put $k=\Q(\alpha)$. Then $k$ is totally real. Its ring of integers is $\mathcal O_k=\{a\in k:a\text{ is integral over }\Z\}$. The polynomial is not assumed monic, and no equality $\mathcal O_k=\Z[\alpha]$ is assumed.

The ring $\mathcal O_k$ is a principal ideal domain if every ideal is principal, equivalently its ideal class number is $1$. It is Euclidean for the standard absolute field norm if
\[
 \forall a,b\in\mathcal O_k\ (b\ne0)\quad
 \exists q\in\mathcal O_k:\quad
 |N_{k/\Q}(a-bq)|<|N_{k/\Q}(b)|.
\]
This asks for the norm itself as a Euclidean function, not merely some Euclidean algorithm.

Long's item 7.1 asks how often each property occurs, and specifically whether the PID property occurs infinitely often. The meaningful field-theoretic infinitude question is whether infinitely many pairwise nonisomorphic totally real fields have class number $1$; this intended reading must be distinguished from infinitely many polynomial presentations of the same field. No fixed degree or ordering for frequency is specified. Counting fields by discriminant, counting polynomials by height, and varying the degree are different problems, so no particular density is asserted.

The source notes that standard-norm Euclideanity occurs only finitely often among quadratic fields. This is retained as supplied background, not a conclusion about all degrees.
\subsection{Short English statement}
Among totally real number fields described by irreducible real-rooted integral polynomials, how frequent are class number one and norm-Euclidean integer rings? Does class number one occur for infinitely many distinct fields?
\subsection{Sources}
\begin{itemize}
\item Local registry: \url{data/unsolvedmath/problems.json}, permanent UnsolvedMath ID $10900049$, exact source code \texttt{catalogue entry}, attributed to Long.
\item Delp, Hoffoss and Manning, \emph{Problems In Groups, Geometry, and Three-Manifolds} (2015), subsection 7.1, PDF page 5: \url{https://arxiv.org/pdf/1512.04620#page=5}; \url{https://arxiv.org/abs/1512.04620}. Supplied access date: 2026-07-29; source-tex extraction checked against rendered PDF text by the source record.
\item The quadratic norm-Euclidean finiteness note and the registry's modern class-number-one status assessment have no separate primary citations in the supplied record.
\end{itemize}
\end{document}
